% ---------------------------------------------------------------------------
\subsection{Cross-Stack Identifiability Audit: GRPO and PPO Saturation Fits on the Same Rollouts}
\label{sec:scaling-law-iter29}
% ---------------------------------------------------------------------------

\paragraph{Motivation.}
Iter~25 audited the canonical saturation model
$R(t)=R_{\max}\bigl(1-e^{-\lambda t}\bigr)$ on the five frontier-scale
traces (Qwen3.5-4B, Qwen3-8B, Llama-3.1-8B-Instruct, DeepSeek-V3.1,
Nemotron-120B) and concluded that the model is \emph{not identifiable}
on those data: 4/5 traces had $\lambda$ pinned at the upper bound
$\lambda=10$ (Nemotron-120B is the exception at $\lambda\!\approx\!0.99$),
0/5 preferred saturation over a constant in noise-aware
AICc, and 0/5 were identifiable~\citep{tinker-rl-lab-iter25}.
That audit left open whether the failure is intrinsic to the model or
an artefact of the 20--30 step frontier horizons (where the noise floor
is comparable to the signal). Iter~29 closes that gap by re-running the
same identifiability battery on the longer-horizon same-stack traces
(40 steps, $n_{\mathrm{eff}}=128$, GSM8K, Qwen2.5-0.5B), and---crucially---
on \emph{both} algorithm stacks simultaneously. The frontier synthesis
Round~1 licenses an \emph{Estimator-Equivalence Principle} (EEP): with
rollout batch, KL, clipping, masking, optimiser and reward parser fixed,
PPO and GRPO should be performance-equivalent. If EEP holds, the
saturation-fit identifiability profile of the two stacks should be
indistinguishable.

\paragraph{Method.}
The driver \texttt{scripts/scaling\_law\_iter29.py} consumes the ten
same-stack runs from \texttt{samestack\_ppo\_grpo.json} (5 GRPO + 5 PPO
seeds, identical everything else) and applies the iter~25 noise-aware
diagnostics per trace:
(a) saturation fit with $\lambda$ bounded in $[10^{-3},\,10]$; (b)
effective binomial $n$ inferred from the GCD of reward denominators;
(c) AICc model selection over $\{\text{constant}, \text{linear},
\text{saturation}\}$; (d) parametric bootstrap ($B=400$) on
$\mathrm{Binomial}(n_{\mathrm{eff}},\,p_t)$ for the lambda CI; (e)
detectability $T^{\star}=(z_{1-\alpha/2}+z_{\mathrm{power}})^{2}$
benchmark against the actual trace length.

Three EEP-falsifiable predictions were recorded in the internal iteration ledger (not an immutable external preregistration):
\begin{description}
\item[E1.] lambda-at-bound rate: PPO $=$ GRPO (Fisher exact, 2x2).
\item[E2.] AICc-best saturation rate: PPO $=$ GRPO (Fisher exact, 2x2).
\item[E3.] Heldout-residual explained by the saturation fit:
  PPO $=$ GRPO (Welch $t$, two-sided).
\end{description}
A non-rejected E1 $\wedge$ E2 $\wedge$ E3 would sustain EEP at the
\emph{shape} level; rejection of any cell is a controlled
\emph{partial} falsification.

\paragraph{Results.}

\begin{table}[t]
\centering\small
\begin{tabular}{lcccccc}
\toprule
Algo & $n$ & lam-at-bound & AICc-best sat & AICc-best lin & AICc-best const & $\hat\lambda$ (med.) \\
\midrule
GRPO & 5 & 0/5 & \textbf{5/5} & 0/5 & 0/5 & 0.250 \\
PPO  & 5 & 0/5 & 0/5 & \textbf{5/5} & 0/5 & 0.078 \\
\bottomrule
\end{tabular}
\caption{Cross-stack identifiability battery on the 40-step same-stack
  traces. Lambda is identifiable in both stacks (0/5 hit the upper
  bound, vs.\ 5/5 in iter~25's 20--30 step frontier data). However the
  AICc-best model \emph{diverges}: GRPO's running mean reward is best
  described by an exponential saturation curve (5/5), PPO's by a linear
  trend (5/5). Fisher exact on the AICc-best-saturation rate gives
  $p=0.0079$ (Table~\ref{tab:iter29-eep}).}
\label{tab:iter29-cross-stack}
\end{table}

\begin{table}[t]
\centering\small
\begin{tabular}{llccc}
\toprule
Prediction & GRPO & PPO & $p$ & EEP \\
\midrule
E1 lambda-at-bound rate       & 0/5        & 0/5        & 1.0000  & sustained \\
E2 AICc-best sat rate         & 5/5        & 0/5        & \textbf{0.0079} & \textbf{falsified} \\
E3 heldout-resid from sat fit & 0.1023     & 0.3901     & 0.0537  & sustained \\
F1 bootstrap CI excludes bound & 5/5       & 4/5        & n/a     & sustained \\
F2 paired heldout accuracy     & 0.99$\pm$0.0035 & 0.992$\pm$0.003 & 0.3739 & sustained \\
\bottomrule
\end{tabular}
\caption{EEP-falsification battery on the same-stack traces.
  E2 is decisively falsified: GRPO and PPO prefer different AICc-best
  functional forms (saturation vs.\ linear) on the same rollout batches.
  E1 and F1 are sustained: both stacks have identifiable $\lambda$ (0/5
  hit the bound), and the bootstrap CI excludes the bound in 5/5 GRPO
  and 4/5 PPO traces. E3 borderline ($p=0.054$): PPO's saturation fit
  leaves more heldout residual than GRPO's, consistent with E2's
  AICc verdict. F2 (final accuracy) is sustained, matching iter~9.}
\label{tab:iter29-eep}
\end{table}

\paragraph{Interpretation.}
The strong Estimator-Equivalence Principle \emph{at the level of
functional-form identifiability} is \textbf{falsified} on this data:
GRPO's running reward is well-described by
$R(t)\approx R_{\max}(1-e^{-\lambda t})$ with $\hat\lambda\approx 0.25$
and $\hat R_{\max}\approx 0.89$ (a clean $t_{80}\approx 6.4$ steps);
PPO's running reward is \emph{not}---its preferred model is linear with
a much smaller effective rate ($\hat\lambda\approx 0.08$ if one forces
the saturation form, vs.\ the GRPO estimate of $0.25$). This is
\emph{not} a final-accuracy effect: the two stacks are statistically
indistinguishable on heldout accuracy ($p=0.37$), and the bootstrap CI
on $\lambda$ is well-defined in both stacks (0/5 hit the bound, vs.\
5/5 in the 20--30 step frontier data).

The decoupling---identical heldout accuracy, divergent
intermediate-curve shape---is the sharpest form of the critic
degeneracy pattern flagged in the frontier synthesis: PPO's value
head appears to flatten the intermediate gradient signal into a
near-linear drift rather than a saturating exponential. Practically,
this means that \emph{rules-of-thumb derived from GRPO saturation fits
do not transfer to PPO} even when everything else in the stack is held
fixed. The two algorithms reach the same accuracy, but they get there
via different intermediate trajectories---and the saturation model
$R(t)=R_{\max}(1-e^{-\lambda t})$ is only the right descriptive model
for one of them.

\paragraph{Limitations.}
The same-stack data are at a single model size (Qwen2.5-0.5B), single
task (GSM8K), and 40-step horizon. Whether the same AICc divergence
holds at 8B+ with the frontier horizon needs a Tinker reproduction
(cost-flagged in \texttt{FRONTIER\_INSIGHTS.md}). The
``saturation-supported'' criterion in iter~25 (AICc-best $=$ sat AND
CI excludes bound) is satisfied in 5/5 GRPO traces and 0/5 PPO traces;
iter~25's stricter formulation would also falsify EEP here.

\paragraph{Artifacts.}
Driver: \texttt{scripts/scaling\_law\_iter29.py}.
Per-trace: \texttt{scaling\_law\_iter29\_identifiability.tsv},
\texttt{scaling\_law\_iter29\_bootstrap.tsv}.
Stack rollup: \texttt{scaling\_law\_iter29\_summary.tsv}.
EEP battery: \texttt{scaling\_law\_iter29\_stack\_compare.tsv}.
Figure: \texttt{figures/scaling\_law\_iter29.pdf}.